\usepackage[utf8x]{inputenc}
\usepackage{ucs}
\usepackage{amsmath}
\usepackage{amsfonts}
\usepackage{amssymb}
\usepackage{graphicx}
\usepackage{appendixnumberbeamer}
\usepackage{url}

\title{GPL, MIT, ..., WTF?}
%\subtitle{}
\author{Urs Fässler}
\date{26.11.2015}
\institute
{
  LinuxDay Dornbirn
}
%\subject{}

\newcounter{todonr}
\newcommand{\todo}[1]{{\color{red}TODO \arabic{todonr}: #1\stepcounter{todonr}}}

\newcommand{\hnote}[1]{\only<handout>{\footnote{#1}}}
\newcommand{\hcite}[1]{\only<handout>{\cite{#1}}}
\bibliographystyle{plainurl}  % urlbst
\setbeamertemplate{bibliography item}{[\theenumiv]}

\usetheme{Luebeck}

% remove navigation bars
\beamertemplatenavigationsymbolsempty

% remove section information
\setbeamertemplate{headline}{}

\begin{document}

\frame{\titlepage}

%\section{Einführung}
\section{Intro}
\note
{
  \begin{itemize}
    \item Teaser: "Kunstwerk erstellen"
  \end{itemize}
}

\section{Thema}
\note
{
  \begin{itemize}
    \item Eine Einführung in die Lizenzen der Freien Software (Open Source)
  \end{itemize}
}

\section{Behauptung}
\note
{
  \begin{itemize}
    \item Korrekte Verwendung Freier Software ist nicht schwer
  \end{itemize}
}

\section{Vorstellung}
\begin{frame}
  \begin{itemize}
    \item Urs Fässler
    \item Senior Embedded Software Entwickler @ bbv Software Service
    \item Erstellen von Systemen auf Basis GNU/Linux
    \item Entwickeln von Software
    \item Etwas involviert in Patentantrag
    \item Aktiv in der FSFE Lokalgruppe Zürich
  \end{itemize}
\end{frame}

\begin{frame}
  \begin{itemize}
    \item Haftungsausschluss
  \end{itemize}
\end{frame}

\section{Einführung}

\begin{frame}{Freie Software}
  4 Freiheiten
\end{frame}

\begin{frame}{Open Source}
\end{frame}

\begin{frame}{FLOSS}
\end{frame}

\begin{frame}{Proprietär versus Kommerziell}
  \begin{center}
    \begin{tabular}{|c|c|c|}
    \hline 
     & proprietär &  frei \\ 
    \hline 
    kommerziell & Windows, Mac OS & Red Hat Enterprise Linux (RHEL) \\ 
    \hline 
    gratis & Freeware & Debian \\ 
    \hline 
    \end{tabular} 
  \end{center}
\end{frame}

\begin{frame}
\begin{center}
Geistiges Eigentum - IP
\end{center}
\end{frame}
\note
{
(IP, Intellectual Property, Geistiges Eigentum) \url{https://de.wikipedia.org/wiki/Geistiges\_Eigentum\#Systematik}
  \begin{itemize}
    \item Arten von Recht, was schützen sie
  \end{itemize}
}

\begin{frame}[fragile]{IP - Markenrecht}
  \todo{Beispiel Iceweasel}
\end{frame}
\note
{
  \begin{itemize}
    \item Markenrecht: \url{https://de.wikipedia.org/wiki/Marke_(Recht)}
    \item Firefox und Debian (Iceweasel): \url{https://lwn.net/Articles/676799/}
  \end{itemize}
}

\begin{frame}{IP - Patent}
  \todo{Beispiel mp3}\\
  \todo{Private versus Geschäftlich}\\
  \todo{Softwarepatente}
\end{frame}
\note
{
  \begin{itemize}
    \item Patent: \url{https://de.wikipedia.org/wiki/Patent}
    \item Defensivpublikation
  \end{itemize}
}

\begin{frame}{IP - Urheberrecht}
  \todo{Beispiel Urheberrecht}\\
  \todo{\textcopyright}
\end{frame}
\note
{
  \begin{itemize}
    \item Urheberrecht: \url{https://de.wikipedia.org/wiki/Urheberrecht}
    \item gilt automatisch: \url{https://anwalt-im-netz.de/urheberrecht/urheberrecht-faq.html\#wann}
    \item Copyrightzeichen ist nicht nötig: \url{https://de.wikipedia.org/wiki/Copyrightzeichen}
    \item Nicht verzichtbar (kein public domain, Gemeinfreiheit: \url{https://de.wikipedia.org/wiki/Gemeinfreiheit\#Entlassung\_in\_die\_Gemeinfreiheit})
  \end{itemize}
}

\section{Inhalt}

\subsection{Lizenzen}
\begin{frame}
  Lizenzen
\end{frame}

\begin{frame}
  \todo{Wissen die Teilnehmer was kompilieren und linken ist?}
  \begin{itemize}
    \item Überblick: \url{https://tldrlegal.com/}
    \item Copyleft
  \end{itemize}
\end{frame}

\begin{frame}{Copyleft}
  \begin{itemize}
    \item viral
  \end{itemize}
\end{frame}
\note
{
  \begin{itemize}
    \item Copyleft, viral: https://de.wikipedia.org/wiki/Copyleft
  \end{itemize}
}

\begin{frame}{Übersicht}
  \begin{description}
    \item[Public Domain] The Unlicense, Beerware, WTFPL
    \item[Nachgiebig] MIT, BSD, Apache, X11
    \item[Copyleft] GPL, LGPL, AGPL, MPL
  \end{description}
\end{frame}

\begin{frame}{Public Domain Lizenzen}
  \begin{itemize}
    \item Nicht möglich, nicht gültig
    \item sollten nicht benutzt werden
    \item Verwende CC-0
  \end{itemize}
\end{frame}
\note
{
	\begin{itemize}
		\item WTFPL (\url{https://en.wikipedia.org/wiki/WTFPL})
		\item Beerware
		\item The Unlicense (\url{http://programmers.stackexchange.com/questions/147111/what-is-wrong-with-the-unlicense\#147120})
	\end{itemize}
}

\begin{frame}{Nachgiebige Lizenzen}
  \begin{itemize}
    \item Darf
    \begin{itemize}
      \item Ist generell Freie Software
      \item Kommerzielle Verwendung
      \item Proprietäre Verwendung
    \end{itemize}
    \item Muss
    \begin{itemize}
      \item Copyright angeben
      \item Lizenz mitliefern
    \end{itemize}
  \end{itemize}
\end{frame}

\begin{frame}{Copyleft Lizenzen}
  \begin{itemize}
    \item Darf
    \begin{itemize}
      \item Ist generell Freie Software
      \item Kommerzielle Verwendung
    \end{itemize}
    \item Darf nicht
    \begin{itemize}
      \item Proprietäre Verwendung
    \end{itemize}
    \item Muss
    \begin{itemize}
      \item Code veröffentlichen
      \item Copyright angeben
      \item Lizenz mitliefern
    \end{itemize}
  \end{itemize}
\end{frame}

\begin{frame}{Werke: Creative Commons}
	Optionen:\\
	\begin{tabular}{|l|l|l|}
		\hline
		BY	& Attribution		& Frei\\
		\hline
		SA	& Share Alike		& Frei\\
		\hline
		NC	& Non-Commercial	& Unfrei\\
		\hline
		ND	& No Derivatives	& Unfrei\\
		\hline
	\end{tabular} 
	
	Verwende:\\
	\begin{tabular}{|l|l|l|}
		\hline
		CC-0		& kein Copyright wenn möglich\\
		\hline
		CC-BY		& Namensnennung\\
		\hline
		CC-BY-SA	& Namensnennung, Viral\\
		\hline
	\end{tabular} 
\end{frame}

\subsection{Handhabung der Lizenzen}

\begin{frame}{Allgemein}
	\begin{itemize}
		\item Eventuell hat man durch Vertrag die Rechte abgetreten (z.B. durch Arbeitsvertrag)
		\item Rechteinhaber darf mit Code machen was er will
		\begin{itemize}
			\item Mehrfachlizenzierung
			\item verschiedene FLOSS Lizenzen
			\item proprietäre Verwendung
		\end{itemize}
	\end{itemize}
\end{frame}

\begin{frame}{Auswahl einer Lizenz}
	\begin{itemize}
		\item Wieso will man den Code Veröffentlichen?
		\item Was will man schützen?
		\item Was will man erlauben?
		\item Was will man verhindern? (Gewissen Einsatz verhindern?)
		\item Webseite: \url{http://choosealicense.com/}
	\end{itemize}
	\todo{verschiedene Szenarien auflisten, Beispiele}
\end{frame}

\begin{frame}{Deklaration der Lizenz}
	\begin{itemize}
		\item im Projekt (COPYING)
		\item in den Files (SPDX)
		\item in dem Binary (Inkscape)
		\item Keine Lizenz: Code darf nicht benutzt werden
	\end{itemize}
\end{frame}

\begin{frame}{Veröffentlichen der Sourcen}
	\begin{itemize}
		\item öffentliches Repository
		\item Download
		\item Auf Anfrage
	\end{itemize}
\end{frame}

\begin{frame}{Verwendung mit proprietären Code}
	\begin{itemize}
		\item proprietäre Applikation auf freiem OS (z.B. Debian GNU/Linux) kein Problem
		\item Verwenden in proprietärem Code: nur nachgiebige Lizenzen (MIT)
		\item Linken nur gegen Libraries mit nachgiebige Lizenzen oder LGPL
		\item Code unter Copyleft darf nicht in proprietärer Applikation verwendet werden (ausser wenn ausschliesslich intern gebraucht)
		\item Lizenz Kompatibilität beachten beim mischen von Code unter verschiedenen Copyleft Lizenzen
	\end{itemize}
\end{frame}

\begin{frame}{Verwendung von Code unter Copyleft Lizenz}
	\begin{itemize}
		\item Lizenz Kompatibilität beachten beim mischen von Code unter verschiedenen Copyleft Lizenzen
		\item Contributoren sollten Patches unter gleicher Lizenz einreichen
		\item Contributions schwierig wenn Dual-Lizenziert
		\item Contributor License Agreement
	\end{itemize}
\end{frame}

\begin{frame}{Verteidigung des eigenen Codes}
	\begin{itemize}
		\item Dialog suchen
		\item gpl-violations.org \url{https://en.wikipedia.org/wiki/Gpl-violations.org}
		\begin{itemize}
			\item Verschiedene Erfolge, dadurch Gültigkeit der GPL Bewiesen
		\end{itemize}
	\end{itemize}
\end{frame}

\begin{frame}{Best Practises in der FOSS Community}
	\begin{itemize}
		\item release early
		\item release often
		\item offener Entwicklungsprozess
		\item einfach Kontaktierbar
		\item nett sein
	\end{itemize}
\end{frame}

\section{Message}
\begin{frame}{Message}
	\begin{itemize}
		\item Korrekter Einsatz Freier Software ist nicht schwer
		\item Korrekte Deklaration Freier Software macht die Verwendung einfacher
	\end{itemize}
\end{frame}

\only<handout>
{
\appendix

\section{Backup}


\section{Outro}
\begin{frame}{Danke}
	\visible<2->{
		\begin{center}
			\includegraphics[height=2cm]{res/Question-Girl.pdf}\hcite{questionGirl}
			\includegraphics[height=2cm]{res/Question-Guy.pdf}\hcite{questionGuy}
		\end{center}
	}
%	\begin{itemize}
%		\item Slides auf \url{http://www.bitzgi.ch/presentations}
%		\item \href{https://github.com/ursfassler/embedded-gnu-linux-basics}{/embedded-gnu-linux-basics} Sourcen des Vortrags
%	\end{itemize}
	\begin{center}
		\includegraphics[width=2cm]{res/by-sa.pdf}
		\hspace{0.5cm}
		Urs Fässler
		\hspace{0.25cm}
		\includegraphics[width=2.5cm]{res/mail.pdf}
	\end{center}
\end{frame}


\begin{frame}{GPL-3}
  \todo{Übersichtlicher machen}\\
  \todo{hier oder bei der Verwendung genauer erklären?}\\
  \begin{tabular}{|c|c|c|}
    \hline 
    Can & Cannot & Must \\ 
    \hline 
    Commercial Use & Sublicense & Include Original \\ 
    \hline 
    Modify & Hold Liable & State Changes \\ 
    \hline 
    Distribute &  & Disclose Source \\ 
    \hline 
    Place Warranty &  & Include License \\ 
    \hline 
    Use Patent Claims &  & Include Copyright \\ 
    \hline 
    &  & Include Install Instructions \\ 
    \hline 
  \end{tabular} 
\end{frame}
\note
{
  \url{https://tldrlegal.com/license/gnu-general-public-license-v3-(gpl-3)}
}


\section{Literatur \& Bild-Nachweise}
\begin{frame}[allowframebreaks]{Literature}
	\tiny{
		\bibliography{literature}
	}
\end{frame}
}

\end{document}
